\chapter{The master inequality: extension to smooth functions}

The aim of this chapter is to prove the Master inequality II, as
done in \cite[Section 7]{chen2025}.
This will require to prove some results related to approximations
of smooth functions by Chebyshev polynomials on a compact set.

Throughout the rest of this blueprint,
we will denote as $\tor = \R / 2 \pi\Z$ the additive circle of length $2 \pi$.

\section{Chebyshev polynomials}

For $j \geq 0$, we denote as $T_j$ the Chebyshev polynomial of the first kind of degree $j$,
defined by the relation $T_j(\cos \theta) = \cos(j \theta)$ for $\theta \in \R$.

Let us fix a real number $K > 0$ and an integer $q \geq 0$ for this section.
We define coefficients associated to a polynomial $h$ on the set $[-K,K]$ in the following
fashion.

\begin{lemma}
  \label{lem:chebyshev_decomp}
  Let $h \in \mathcal{P}_q$.
  There exists a unique family of real numbers $(a_j)_{0 \leq j \leq q}$ such that,
  for all $x \in \R$,
  \begin{equation*}
    h(Kx) = \sum_{j=0}^q a_j T_j (x).
  \end{equation*}
\end{lemma}

\begin{proof}
  This follows from the fact that $h(Kx) \in \mathcal{P}_q$ and $(T_j)_{0 \leq j \leq q}$ is a basis of $\mathcal{P}_q$,
  since for all $0 \leq j \leq q$, $\deg T_j = j$.
\end{proof}

It will be useful to view the coefficients $(a_j)_{0 \leq j \leq q}$ as trigonometric
Fourier coefficients, which we do by introducing the following function $f$.

\begin{lemma}
  \label{lem:f_def}
  \uses{lem:chebyshev_decomp}
  Let $h \in \mathcal{P}_q$ and $(a_j)_{0 \leq j \leq q}$ be the coefficients from Lemma
  \ref{lem:chebyshev_decomp}.
  Let $f : \tor \rightarrow \R$ be the function defined by $f(\theta) = h(K \cos \theta)$.
  Then, $f$ is an even trigonometric polynomial equal to
  \begin{equation*}
    f(\theta)
    =
      \sum_{j=0}^q a_j \cos(j \theta).
 \end{equation*}
\end{lemma}

\begin{proof}
  By definition of $f$ and $(a_j)_{0 \leq j \leq q}$,
  for all $\theta \in \tor$,
  \begin{equation*}
    f(\theta) = h(K \cos \theta) = \sum_{j=0}^q a_j T_j(\cos \theta).
  \end{equation*}
  By definition of $(T_j)_{j \geq 0}$, $T_j(\cos \theta) = \cos (j \theta)$.
  The evenness comes from the fact that $\cos$ is even.
\end{proof}

We will now provide some bounds on the coefficients $(a_j)_{0 \leq j \leq q}$
in terms of $L^\beta$-bounds on derivatives of $f$.
Recall that, for a function $f \in L^1(\tor)$ and an integer $j \in \Z$,
the $j$-th Fourier coefficient is defined as
\begin{equation*}
  c_j = \frac{1}{2 \pi} \int_0^{2 \pi} f(\theta) e^{-i j \theta} \d \theta.
\end{equation*}

\begin{lemma}
  \label{lem:coeff_of_f}
  Let $f(\theta) = \sum_{j=0}^q a_j \cos(j \theta)$ be an even trigonometric series.
  Then, the Fourier coefficients of $f$ are exactly, for $j \in \Z$,
  \begin{equation*}
    c_j =
    \begin{cases*}
      \frac{1}{2} a_{|j|} & if $0 < |j| \leq q$ \\
      a_0 & if $j=0$ \\
      0 & if $|j| > q$.
    \end{cases*}
  \end{equation*}
\end{lemma}

\begin{proof}
  This is probably already in Mathlib.
\end{proof}


\begin{lemma}
  \label{lem:bound_a0}
  \uses{lem:f_def,lem:coeff_of_f}
  Let $h \in \mathcal{P}_q$ and $(a_j)_{0 \leq j \leq q}$ be the coefficients from Lemma
  \ref{lem:chebyshev_decomp}.
  Then,
  \begin{equation*}
    |a_0| \leq \| h \|_{C^0[-K,K]}.
  \end{equation*}
\end{lemma}

\begin{proof}
  By Lemma \ref{lem:f_def} and \ref{lem:coeff_of_f}, we have
  \begin{equation*}
    a_0 = c_0 = \frac{1}{2 \pi} \int_0^{2 \pi} f(\theta) \d \theta.
  \end{equation*}
  Hence by the triangle inequality
  \begin{equation*}
    |a_0| \leq \frac{1}{2 \pi} \int_0^{2 \pi} |f(\theta)| \d \theta
    = \frac{1}{2 \pi} \int_0^{2 \pi} |h(K \cos \theta)| \d \theta
    \leq \| h \|_{C^0[-K,K]}.
  \end{equation*}
\end{proof}

The following bound relies on the Hausdorff-Young inequality, which is a classic result.

\begin{theorem}
  \label{thm:hausdorff_young}
  Let $1 < \beta \leq 2$ and $1/\beta_* = 1 - 1/\beta$.
  Let $f \in L^\beta(\tor)$ of Fourier coefficients $(c_j)_{j \in \Z}$.
  Then,
  \begin{equation*}
    \left(\sum_{j \in \Z} |c_j|^{\beta_*} \right)^{1/\beta_*}
    \leq \left(\frac{1}{2 \pi} \int_0^{2 \pi} |f(\theta)|^\beta \d \theta\right)^{1/\beta}.
  \end{equation*}
\end{theorem}

\begin{proof}
  I think this might be in Mathlib? If not this is a classic result and should be added.
  This is proven by interpolation between Parseval and the $L^\infty$-$L^1$ trivial estimate,
  using Riesz-Thorin.
\end{proof}

We deduce the following lemma.

\begin{lemma}
  \label{lem:Zygmund}
  \uses{thm:hausdorff_young}
  Let $\beta > 1$ and $1/\beta_* = 1-1/\beta$.
  Let $f \in L^1(\tor)$ of Fourier coefficients $(c_j)_{j \in \Z}$.
  Assume $f' \in L^\beta(\tor)$.
  Then,
  \begin{equation*}
    \sum_{j \neq 0} |c_j| \leq 4 \beta_* \|f' \|_{L^\beta(\tor)}.
  \end{equation*}
\end{lemma}

\begin{proof}
  Assume $1 < \beta \leq 2$.
  By the H\"older inequality,
  \begin{equation*}
    \sum_{j \neq 0} |c_j|
    = \sum_{j \neq 0} \frac{1}{|j|} \times |j c_j|
    \leq
    \left(\sum_{j \neq 0} \frac{1}{|j|^\beta} \right)^{1/\beta}
    \left(\sum_{j \neq 0} |j c_j|^{\beta_*} \right)^{1/\beta_*}.
  \end{equation*}
  We shall bound these two sums successively.
  \begin{itemize}
    \item   Since $\beta>1$ the first sum converges. By parity it is equal to twice the sum for $j \geq 1$.
  Since $t \mapsto 1/t^\beta$
  is decreasing on $[1,\infty)$, we can compare the sum and integral which yields
  \begin{equation*}
    \sum_{j =1}^\infty \frac{1}{j^\beta}
    \leq 1 + \int_1^\infty \frac{\d t}{t^\beta}
    = 1 - \left[ \frac{1}{(\beta-1)t^{\beta-1}} \right]_1^\infty
    = 1 + \frac{1}{\beta - 1}
    = \frac{\beta}{\beta - 1} = \frac{1}{1 - 1/\beta} = \beta_*.
  \end{equation*}
  As a conclusion,
  \begin{equation*}
    \left(\sum_{j \neq 0} \frac{1}{|j|^\beta} \right)^{1/\beta}
    \leq 2 \beta_*^{1/\beta} \leq 2 \beta_*
  \end{equation*}
  since $0 \leq 1/\beta \leq 1$ and $\beta_* \geq 1$.
  \item For the second sum, we observe that the Fourier coefficients of $f'$
  are $(i j c_j)_{j \in \Z}$. Since $1 < \beta \leq 2$, we can
  apply Theorem \ref{thm:hausdorff_young} to $f'$, which yields
  \begin{equation*}
    \left(\sum_{j \neq 0} |j c_j|^{\beta_*} \right)^{1/\beta_*}
    \leq
    \left(\frac{1}{2 \pi} \int_0^{2 \pi} |f'(\theta)|^\beta \d \theta\right)^{1/\beta}
    \leq \frac{1}{(2 \pi)^{1/\beta}} \|f'\|_{L^\beta(\tor)}
    \leq \frac{1}{\sqrt{2 \pi}} \|f'\|_{L^\beta(\tor)}
  \end{equation*}
  since $\beta \leq 2$.
  \end{itemize}
  Putting everything together we obtain
  \begin{equation*}
    \sum_{j \neq 0} |c_j|
    \leq \sqrt{\frac{2}{\pi}} \beta_*  \|f'\|_{L^\beta(\tor)}
  \end{equation*}
  which implies the claim in the case $1 < \beta \leq 2$ since $\sqrt{2/\pi} \leq 4$.

  Now assume $\beta \geq 2$.
  By the H\"older inequality, since $\tor$ has finite measure $2 \pi$,
  $L^\beta(\tor) \subseteq L^2(\tor)$ and
  \begin{equation*}
    \|f'\|_{L^2(\tor)}^2
   \leq  \left(\int_0^{2 \pi} 1 \d \theta\right)^{1-\frac{2}{\beta}}
    \left(\int_0^{2 \pi} |f'(\theta)|^{\beta} \d \theta\right)^{\frac{2}{\beta}}
    = (2 \pi)^{1-\frac{2}{\beta}}  \|f'\|_{L^\beta(\tor)}^2
    \leq 2 \pi \|f'\|_{L^\beta(\tor)}^2.
  \end{equation*}
  In particular, we can apply the result for $\beta=2$ (in which case $\beta_*=2$),
  which yields
  \begin{equation*}
    \sum_{j \neq 0} |c_j|
    \leq 2 \sqrt{\frac{2}{\pi}}  \|f'\|_{L^2(\tor)}
    \leq 4 \|f'\|_{L^\beta(\tor)}
    \leq 4 \beta_*  \|f'\|_{L^\beta(\tor)}
  \end{equation*}
  since $\beta_* \geq 1$.
\end{proof}

We now conclude with the main result of this section.

\begin{lemma}
  \label{lem:bound_sum_aj}
  \uses{lem:Zygmund,lem:coeff_of_f,lem:f_def}
  Let $h \in \mathcal{P}_q$, and $f$ and $(a_j)_{0 \leq j \leq q}$ be as in Lemma \ref{lem:f_def}.
  For any $m \geq 0$,
  any $\beta > 1$, if we denote $1/\beta_* := 1 - 1/\beta$, we have
  \begin{equation*}
    \sum_{j=1}^q j^m |a_j| \leq 4 \beta_* \|f^{(m+1)}\|_{L^\beta(\tor)}.
  \end{equation*}
\end{lemma}

\begin{proof}
  First, we apply Lemma \ref{lem:coeff_of_f} to express the Fourier coefficients
  $(c_j)_{j \in \Z}$ in terms of $(a_j)_{0 \leq j \leq q}$.
  Let $m \geq 0$. By differentiation,
  the Fourier coefficients of $f^{(m)}$ satisfy, for
  all $j$, $c_j^{(m)} =
    i^m j^m c_j$.
  In particular,
  \begin{equation*}
    \sum_{j \neq 0} |c_j^{(m)}|
    = \sum_{j \neq 0} |j|^m |c_j|
    = \sum_{j =1}^q j^m |a_j|.
  \end{equation*}
  The conclusion follows by applying Lemma \ref{lem:Zygmund} to the function $f^{(m)}$,
  which is in $L^\beta(\tor)$ as it is a trigonometric polynomial and in particular smooth.
\end{proof}


\section{Compactly supported distributions}

In this section, we will prove a few useful elementary results on distributions.
We denote as $C^\infty(\R)$ the space of smooth functions $\R \rightarrow \R$, and $C^\infty_c(\R)$
those with compact support. For $h \in C^\infty_c(\R)$, we denote as $\supp h$ its support.
For real numbers $a < b$ and $h \in C^\infty(\R)$, we write
\begin{equation*}
  \|h\|_{C^m[a,b]} = \sup_{0 \leq k \leq m} \sup_{a \leq x \leq b} |h^{(k)}(x)|
\end{equation*}
where $h^{(k)}$ is the $k$-th order derivative of $h$.
We denote as $\mathcal{D}'(\R)$ the space of classical distributions on $\R$, i.e. continuous
linear functionals $\nu : C^\infty_c(\R) \rightarrow \R$.
Here continuous means that, for any real number $K \geq 0$,
there exists a constant $C \geq 0$ and an integer $m \geq 0$ (both depending on $K$) such that
\begin{equation*}
  \forall h \in C^\infty_c(\R), \quad
  \supp h \subseteq [-K,K] \Rightarrow  |\nu(h)| \leq C \|h\|_{C^m[-K,K]}.
\end{equation*}

We recall the definition of support, which should be added quite soon to Mathlib.

\begin{definition}
  \label{def:distribution_vanishes}
  Let $A$ be a subset of $\R$. We say $\nu \in \mathcal{D}'(\R)$ \emph{vanishes on} $A$
  if, for all $h \in C_c^\infty(\R)$, if $\supp h \subseteq A$, then $\nu(h)=0$.
\end{definition}

\begin{definition}
  \label{def:support_distribution}
  \uses{def:distribution_vanishes}
  Let $\nu \in \mathcal{D}'(\R)$.
  The \emph{support} of $\nu$ is the intersection of all closed sets $A \subseteq \R$
  such that $\nu$ vanishes on the complement of $A$. We denote it as $\supp \nu$.
\end{definition}

\begin{definition}
  \label{def:compactly_supported_distribution}
  \uses{def:support_distribution}
  A distribution $\nu \in \mathcal{D}'(\R)$ is \emph{a compactly supported distribution}
  if its support is compact.
\end{definition}

\begin{lemma}
  \label{lem:extension_csd}
  \uses{def:compactly_supported_distribution}
  Let $\nu \in \mathcal{D}(\R)$ be a compactly supported distribution.
  Then there is a unique continuous
  linear extension $\nu : C^\infty(\R) \rightarrow \R$ of $\nu$. Furthermore,
  there exist real numbers $K \geq 0$, $C>0$ and an integer $m \geq 0$
  such that
  \begin{equation*}
    \forall h \in C^\infty(\R), \quad |\nu(h)| \leq C \|h\|_{C^m[-K,K]}.
  \end{equation*}
\end{lemma}

\begin{proof}
  In Rudin Functional Analysis according to Wikipedia.
\end{proof}

This allows us to make sense of $\nu(h)$ for $\nu$ a compactly supported distribution
and $h$ a polynomial, an essential component of the proof.

\begin{definition}
  \label{def:radius_csd}
  \uses{lem:extension_csd}
  Let $\nu \in \mathcal{D}(\R)$ be a compactly supported distribution.
  We define $$\rho(\nu) := \limsup_{p \to \infty} |\nu(x^p)|^{1/p} .$$
\end{definition}

\begin{lemma}
  \label{lem:radius_finite}
  \uses{def:radius_csd}
  Let $\nu \in \mathcal{D}(\R)$ be a compactly supported distribution.
  Then $\rho(\nu) < \infty$.
\end{lemma}

\begin{proof}
  By Lemma \ref{lem:extension_csd}, there exists $K \geq 0$, $C \geq 0$ and an integer $m$
  such that, for all $h \in C^\infty(\R)$, $|\nu(h)| \leq C \|h\|_{C^m[-K,K]}$.
Let $p \geq 0$ be an integer.
Then
$$|\nu(x^p)| \leq C \|x^p\|_{C^m[-K,K]}
= C \sup_{0 \leq k \leq m} \sup_{|x| \leq K} |p(p-1) \ldots (p-m+1) x^{p-m}|
\leq C p^m (K+1)^p$$
and hence $(|\nu(x^p)|^{1/p})_{p \geq 0}$ is bounded.
\end{proof}

\begin{lemma}
  \label{lem:support_limsup}
  \uses{lem:radius_finite}
  Let $\nu \in \mathcal{D}(\R)$ be a compactly supported distribution.
  Let $\nu \in \mathcal{D}(\R)$ be a compactly supported distribution.
  Then $\supp \nu \subseteq [-\rho(\nu), \rho(\nu)]$.
\end{lemma}

\begin{proof}
  Lemma 4.9 in \cite{chen2025}
\end{proof}
